\documentclass[10pt,a4paper]{amsart}
\usepackage[margin=19mm]{geometry}
\usepackage{amsmath,amssymb,mathtools}
\usepackage[scr=rsfs,cal=euler]{mathalfa}
\usepackage{xfrac}
\usepackage[unicode,hidelinks,bookmarks=false]{hyperref}
\newcommand{\ie}{i.e.\ }
\newcommand{\cf}{cf.\ }
\newcommand{\resp}{respectively\ }
\newcommand{\Subsection}{\subsection}
\providecommand{\nobreakdash}{\nobreak\hbox{-}}
\setlength{\emergencystretch}{2em}
\multlinegap0pt
\usepackage{bm}
\usepackage{bbm}
\usepackage{dsfont}
\usepackage{xspace}
\usepackage{cancel}
\usepackage{mathtools}
\usepackage{amsthm}
\makeatletter
\@ifundefined{lemma}{\newtheorem{lemma}{Lemma}}{}
\makeatother
\DeclareMathOperator{\divv}{\mathrm{div}}
\newcommand{\R}{\mathbb{R}}
\newcommand{\abs}[1]{\left \lvert #1 \right \rvert}
\title[Elliptic unique continuation below the Lipschitz threshold]{Elliptic unique continuation below the Lipschitz threshold}
\author{Cole Jeznach}
\date{Source: arXiv:2507.23614v2 (CC BY 4.0)}
\begin{document}
\maketitle

\noindent\emph{Source and licence.} Elliptic unique continuation below the Lipschitz threshold. Version arXiv:2507.23614v2 (CC BY 4.0), distributed under the Creative Commons Attribution 4.0 International licence, as recorded in the arXiv licence metadata for this exact version and shown on the abstract page https://arxiv.org/abs/2507.23614v2. Full rights notice: licenses/arxiv-2507.23614-CC-BY-4.0.txt.\par\smallskip

\noindent\textit{Editorial scope.} Counted unit: the Lemma whose source label is \texttt{lem:almost\_monoton} in the article (source0.tex lines 297--311, source label \texttt{lem:almost\_monoton}), delivered together with the article's own complete original proof of it (source0.tex lines 312--381, one \texttt{proof} environment of 70 lines). No mathematics is written, repaired or re-derived: statement and proof are the article's own text line for line. Required context retained uncounted: eqn:soln (lines 281--283). Every definition, display and label of those lines is retained unchanged. Omitted from this package and not redistributed: 1--280, 284--296, 382--633. Nothing inside a retained excerpt is omitted or altered: every retained body is delivered exactly as the article's lines are written, so no omission receipt of any kind accompanies this package. Citations: the retained excerpts carry no citation command at all, so no bibliography entry of the article is transcribed and no reference is redistributed. Reference adaptation: none was needed and none was made; every reference inside the delivered unit resolves inside it, so no unresolved reference or citation is produced. Printed slips kept literal: no printed word, symbol, hypothesis or number of the counted statement or of its proof was altered. Layout and class: the article's own class is \texttt{amsart} with packages including amsmath, amssymb, amsthm, hyperref, esint, enumerate, enumitem, mathtools, comment, babel, xcolor, geometry; the wrapper is the lane's pinned \texttt{amsart} editorial wrapper at 10pt on A4, which restates the theorem-like environments the retained text needs and adds the article's own macro definitions that the retained text needs (\texttt{\textbackslash R}, \texttt{\textbackslash abs}, \texttt{\textbackslash divv}), with extra wrapper packages (bm, bbm, dsfont, xspace, cancel, mathtools). The delivered numbering is the unit's own. Assets: no retained span contains a figure environment, a graphics-inclusion command, a PDF-page inclusion, a table or a TikZ picture; no image file is copied into this package, the package declares no asset dependency and no image file is redistributed with it. Licence: version arXiv:2507.23614v2 of this article is distributed under the Creative Commons Attribution 4.0 International licence, as recorded in the arXiv licence metadata for this exact version and shown on the abstract page https://arxiv.org/abs/2507.23614v2. A full rights notice is delivered at licenses/arxiv-2507.23614-CC-BY-4.0.txt. Characters: every retained excerpt is pure ASCII, so the delivered engine, XeLaTeX with its default Latin Modern text font, reports no missing character.
	\begin{align}
		-\divv(A \nabla u) & = 0  \text{ in } B_2 \subset \R^n, \label{eqn:soln}
	\end{align}

	\begin{lemma}\label{lem:almost_monoton}
		Let $u$ be a solution of \eqref{eqn:soln} in $B_R$, and suppose that $A$ is symmetric and uniformly elliptic with constant $\lambda >0$, and $A$ satisfies the smoothness estimates 
		\begin{align*}
			\abs{\nabla A(x)} \le M, \qquad \abs{A(x) - I } \le \delta, \, \, x \in B_R, 
		\end{align*}
		for some $ M, \delta >0$. Then we have the estimate 
		\begin{align*}
			\dfrac{d}{dr}  \log \left(   N_u^A(r) \right) \ge - C_1 ( M + \delta/r), \, \, r \in (0,R),
		\end{align*}
		for some constant $C_1 >0$ depending just on the dimension $n$ and $\lambda$. In addition, we have the estimate
		\begin{align*}
			\dfrac{d}{dr} \log( r^{-n+1} H_u^A(r)) & =  2 N_u^A(r)/r + e(r), \, \, r \in (0, R),
		\end{align*}
		with error $\abs{e(r)} \le C_1 ( M + \delta/r)$. 
	\end{lemma}

	\begin{proof}
		Recall from the definition of $N_u^A$, we have that $\mu(x) \equiv \langle A(x)x/\abs{x}, x/\abs{x}\rangle$. Define the vector field $\beta(x) \coloneqq A(x)x/\mu(x)$. Since the unit outer normal $\nu(x)$ to $\partial B_r$ is given by $\nu(x) = x/\abs{x}$, then we have by definition of $\mu$ that
		\begin{align}
			A(x) \nu(x) \equiv \mu(x) \beta(x)/r, \, \,  \beta(x) \cdot \nu(x) = r, \, \, x \in \partial B_r,  \label{eqn:beta_norm}
		\end{align}
		Straight-forward arguments using the ellipticity and smoothness conditions on $A$ give us the following estimates on $\mu$ and $\beta$:
		\begin{align}
			& \mu(x)  \simeq 1, \, \abs{\nabla \mu(x)} \lesssim M + \delta/\abs{x}, \, \abs{\mu(x)  -1} \lesssim \delta, \label{eqn:mu_est} \\
			& \abs{\beta(x) - x} \lesssim \delta \abs{x}, \, \abs{D\beta(x) - I} \lesssim  M \abs{x}  + \delta, \label{eqn:beta_est}
		\end{align}
		with implicit constants depending just on the dimension and $\lambda$. Of course in order to prove the estimates above, one uses that if $A \equiv I$, then $\mu(x) \equiv 1$ and $\beta(x) \equiv x$. 
		For convenience, write $N(r) = N_u^A(r)$, $D(r) = D_u^A(r)$ and $H(r) = H_u^A(r)$. One computes
		\begin{align}\label{eqn:logN'}
			\dfrac{d}{dr} \log(N(r)) & = \dfrac{1}{r}  + \dfrac{D'(r)}{D(r)} - \dfrac{H'(r)}{H(r)},
		\end{align}
		and we deal with each term separately.
		
		First, as for $D'(r)$ we have from \eqref{eqn:beta_norm} and the divergence theorem that
		\begin{align*}
			D'(r) & = \int_{\partial B_r} \langle A \nabla u, \nabla u \rangle \\
			& = r^{-1} \int_{\partial B_r} \langle A \nabla u, \nabla u \rangle \langle \beta, \nu \rangle  \\
			& = r^{-1} \int_{B_r} \divv( \langle A \nabla u, \nabla u \rangle \beta ) ,
		\end{align*}
		and so using the symmetry of $A$,
		\begin{align*}
			rD'(r) & = \int_{B_r} \divv(\beta) \langle A \nabla u, \nabla u \rangle + \int_{B_r} \sum_{i,j,k} \beta^k (A_{ij})_{x_k} u_{x_i} u_{x_j} + 2 \int_{B_r} \sum_{i,j,k} \beta^k  A_{ij} u_{x_i} u_{x_j x_k} \\
			& \coloneqq T_1 + T_2 + T_3.
		\end{align*}
		By courtesy of \eqref{eqn:beta_est} and the fact that $\mathrm{Trace}(I) = n$, we have
		\begin{align*}
			T_1 = n D(r) + O(M r + \delta) D(r) , \,\, T_2 = O(M r) D(r).
		\end{align*}
		As for $T_3$, we integrate by parts, use that $u$ solves \eqref{eqn:soln}, and again apply \eqref{eqn:beta_norm}, \eqref{eqn:beta_est} to obtain 
		\begin{align*}
			T_3 & = -2 \int_{B_r} \sum_{ijk} (\beta^k)_{x_j} A_{ij} u_{x_i} u_{x_k} + 2 \int_{\partial B_r} \langle \beta, \nabla u \rangle \langle A \nabla u, \nu \rangle \\
			& = -2 D(r) + O(Mr + \delta) D(r) + 2r \int_{\partial B_r} \mu^{-1} \langle A \nabla u, \nu \rangle^2 .
		\end{align*}
		Finally, we use once again that $u$ is a solution to rewrite
		\begin{align}\label{eqn:D}
			D(r) & = \int_{B_r} \divv(u A \nabla u)  = \int_{\partial B_r} u \langle A \nabla u, \nu\rangle. 
		\end{align}
		Combining all of the the previous equalities yields
		\begin{align}\label{eqn:D'}
			D'(r)/D(r) & = \dfrac{n-2}{r} + O(M + \delta/r) + 2 \dfrac{ \int_{\partial B_r} \mu^{-1} \langle A \nabla u, \nu \rangle^2  }{  \int_{\partial B_r} u \langle A \nabla u, \nu \rangle   }.
		\end{align} 
		
		Moving on to $H'(r)$, one easily computes with the estimate on $\mu$ from \eqref{eqn:mu_est} that 
		\begin{align*}
			H'(r) & = \dfrac{n-1}{r} H(r) + \int_{\partial B_r} \partial_\nu ( u^2 \mu) \\
			& = \dfrac{n-1}{r} H(r) + O(M + \delta/r) H(r) + 2\int_{\partial B_r} u \partial_\nu u \mu. 
		\end{align*}
		Next we notice that $A(x) \nu(x) - \mu(x) \nu(x) \equiv r^{-1} \mu(x) (\beta(x) - x)$ is a tangent vector field to $\partial B_r$ which satisfies
		\begin{align*}
			\abs{\divv_{\partial B_r} (A\nu - \mu \nu)} \lesssim M + \delta/r.
		\end{align*}
		Indeed the estimate above again comes from the estimates \eqref{eqn:mu_est}, \eqref{eqn:beta_est}. Thus we can apply the divergence theorem on $\partial B_r$ and use the symmetry of $A$ to obtain
		\begin{align*}
			\abs{ \int_{\partial B_r} 2 u \partial_\nu u \mu  - 2 u \langle A \nabla u, \nu \rangle } & = \abs{ \int_{\partial B_r} \langle \nabla (u^2), \mu \nu - A \nu \rangle   } \\
			& = \abs{ \int_{\partial B_r}  u^2 \divv_{\partial B_r}(\mu \nu - A \nu) } = O(M + \delta/r) H(r).
		\end{align*}
		Moreover, we can use this to rewrite
		\begin{align}\label{eqn:H'}
			H'(r) & = \left( \dfrac{n-1}{r} + O(M + \delta/r) \right) H(r) + 2 \int_{\partial B_r} u \langle A \nabla u, \nu \rangle. 
		\end{align}
		With the estimates \eqref{eqn:D'} and \eqref{eqn:H'}, we substitute into \eqref{eqn:logN'} to obtain
		\begin{align*}
			\dfrac{d}{dr} \left( \log N(r) \right) = 2 \left( \dfrac{\int_{\partial B_r} \mu^{-1} \langle A \nabla u, \nu \rangle^2 }{ \int_{\partial B_r} u \langle A \nabla u, \nu \rangle }  - \dfrac{  \int_{\partial B_r} u \langle A \nabla u, \nu \rangle }{  \int_{\partial B_r} u^2 \mu   }\right) + O(M + \delta/r),
		\end{align*}
		and conclude the proof of the Lemma by Cauchy-Schwarz. 
	\end{proof}

\end{document}
